\documentclass[manuscript,screen]{acmart}

\AtBeginDocument{%
  \providecommand\BibTeX{{%
    Bib\TeX}}}

\settopmatter{printacmref=false}
\renewcommand\footnotetextcopyrightpermission[1]{}
\pagestyle{plain}

\begin{document}

\title[Tracr para Diagnosticar ANC]{Interpretabilidad indirecta ANC a través de Tracr-RASP}

\author{Juan Rodrigo Anabalón Riquelme}
\email{misc@monkeyslab.cl}
\affiliation{%
  \institution{Universidad Nacional de Cuyo}
  \country{Argentina}
}

\author{Juan B. Cabral}
\affiliation{%
  \institution{Universidad Nacional de Córdoba}
  \country{Argentina}
}

\renewcommand{\shortauthors}{Anabalón Riquelme and Cabral}

\begin{abstract}
Las redes de Criptografía Neuronal Adversarial (ANC, \citet{abadi2016learning}) aprenden a
cifrar mensajes de extremo a extremo mediante un juego minimax entre tres redes (Alice, Bob,
Eve), sin que el algoritmo de cifrado resultante sea legible a partir de los pesos entrenados.
En este trabajo proponemos usar un transformer compilado con Tracr
\citep{lindner2023tracr} —que implementa exactamente un cifrado XOR bit a bit a partir de un
programa RASP conocido \citep{weiss2021thinking}— como fuente de un vocabulario de variables
interpretables (bit de clave alineado por posición, bit de salida XOR por posición) que sirve
de hipótesis para sondear linealmente la única capa oculta de una red Alice (un MLP) entrenada
de forma opaca. El sondeo revela una heterogeneidad invisible en la métrica agregada de
coincidencia con XOR: algunos bits de salida son perfectamente legibles de forma lineal,
mientras que otros permanecen a nivel de azar. Mostramos que esta heterogeneidad coincide
exactamente con el colapso de la norma de los pesos de salida correspondientes, consistente
con un óptimo local del objetivo adversarial. Mediante una ablación controlada (barrido de los
pesos de la función de pérdida y del ancho de la capa oculta, 5 semillas por configuración)
mostramos que el fenómeno es principalmente un cuello de botella de capacidad de la
arquitectura original de Abadi y Andersen, y no una limitación fundamental del esquema de
entrenamiento: duplicar el ancho de la capa oculta reduce el colapso de $2.60\pm1.02$ a
$0.40\pm0.49$ posiciones colapsadas (de 8) en promedio, y combinado con un reajuste de los
pesos de la pérdida lo elimina por completo ($\text{xor\_match}=1.000$, 0 posiciones
colapsadas en las 5 semillas). El trabajo ilustra un uso concreto de modelos compilados como
generadores de hipótesis interpretativas para diagnosticar patologías de entrenamiento en
redes opacas.
\end{abstract}

\keywords{interpretabilidad mecanística, Tracr, RASP, criptografía neuronal adversarial, sondeo lineal, transformers compilados}

\maketitle

\section{Introducción}

Abadi y Andersen \citep{abadi2016learning} mostraron que tres redes neuronales entrenadas
adversarialmente —Alice, Bob y Eve— pueden aprender, sin que se les prescriba ningún algoritmo
criptográfico, una forma de cifrado que Bob puede invertir usando una clave compartida y que
Eve, sin acceso a la clave, no puede. El resultado es notable precisamente porque es
\emph{end-to-end}: el cifrado que emerge en los pesos de Alice no es legible directamente, y
los propios autores reportan que, en su formulación original (sin ningún término adicional que
lo empuje hacia una forma particular), el cifrado emergente \emph{no} es simplemente XOR: un
cambio en un bit de clave o de mensaje se difunde de forma continua sobre varios elementos del
texto cifrado, en lugar de afectar un único bit como lo haría un XOR exacto. Para poder comparar directamente contra una referencia con semántica conocida, en este trabajo
agregamos al objetivo de Alice un término auxiliar que explícitamente empuja el cifrado hacia la
forma XOR bit a bit (Sección~\ref{sec:metodo}). Esto no está garantizado por el objetivo
adversarial en sí, sino que es una elección de diseño que hace que XOR sea el blanco declarado
del entrenamiento; pero nada garantiza que ese blanco se alcance de manera completa o uniforme a
lo largo de toda la salida —y es precisamente esa brecha la que investigamos.

Esto plantea un problema de interpretabilidad concreto: ¿cómo se verifica, más allá de una
métrica de precisión agregada, si una red entrenada de forma opaca efectivamente implementa la
operación que se espera que aprenda? Los métodos de interpretabilidad post-hoc habituales
(mapas de saliencia, LIME, SHAP) no fueron diseñados para responder preguntas de esta forma —no
hay, a priori, ninguna referencia contra la cual comparar la explicación producida.

Tracr \citep{lindner2023tracr} ofrece una solución indirecta a este problema: permite escribir
un programa en RASP \citep{weiss2021thinking} —un lenguaje simbólico para razonar sobre
computación en transformers— y compilarlo en un transformer real cuyos pesos implementan
\emph{exactamente} ese programa, con una estructura interna conocida por construcción (variables
del residual stream etiquetadas, patrones de atención predecibles). El objetivo original de
Tracr es servir de ``laboratorio'': un modelo con ground-truth interno conocido sobre el cual
validar técnicas de interpretabilidad antes de aplicarlas a modelos opacos.

En este trabajo llevamos esa idea un paso más allá de su uso habitual. En lugar de usar el
modelo compilado únicamente para validar una técnica de interpretabilidad en abstracto,
usamos el \emph{vocabulario de variables} que expone Tracr para un programa RASP de XOR
—concretamente, la noción de ``bit de clave traído a la posición del mensaje'' y ``bit de
salida XOR en la posición $i$''— como hipótesis concreta para sondear una red Alice
entrenada de forma opaca. Como Alice es un MLP y el modelo compilado es un transformer, la
comparación no puede hacerse a nivel de circuito o de neurona a neurona; el punto de
contraste es conceptual: si Alice aprendió efectivamente una función equivalente a XOR,
debería existir, en algún lugar de su representación interna, una codificación linealmente
legible del bit de salida por posición, análoga a la variable categórica que Tracr expone de
forma explícita.

\subsection{Pregunta de investigación y contribución}

Nuestra pregunta de investigación es: \emph{¿el vocabulario de variables de un modelo
compilado con Tracr puede usarse como generador de hipótesis para un sondeo interpretativo de
una red adversarial entrenada de forma opaca, y ese sondeo revela algo que las métricas
agregadas de desempeño no muestran?} Encontramos que la respuesta es afirmativa, y que además
el sondeo lineal permite diagnosticar mecánicamente una patología de entrenamiento —un colapso
de peso por posición de salida— que resulta ser, mayormente, un problema de capacidad de la
arquitectura y no una limitación intrínseca del esquema ANC. Las contribuciones concretas de
este trabajo son:

\begin{enumerate}
\item Una metodología de sondeo lineal por posición, usando como target las variables
  categóricas expuestas por un programa RASP compilado, aplicada a la única capa oculta de una
  red Alice (Sección~\ref{sec:metodo}).
\item La identificación de una heterogeneidad no uniforme en la convergencia de Alice hacia
  XOR, invisible en la métrica agregada usual, y su explicación mecánica en términos de la
  norma de los pesos de salida por posición (Sección~\ref{sec:resultados}).
\item Una ablación controlada que aísla el efecto de los pesos de la función de pérdida
  adversarial (\(\alpha\), \(\beta\)) del efecto de la capacidad de la capa oculta, mostrando
  que ambos contribuyen pero que la capacidad domina el fenómeno (Sección~\ref{sec:resultados}).
\end{enumerate}

\section{Trabajo relacionado}

\textbf{Criptografía neuronal adversarial.} \citet{abadi2016learning} introdujeron el esquema
Alice--Bob--Eve entrenado mediante un objetivo minimax en el que Alice minimiza el error de
reconstrucción de Bob y es recompensada cuando el error de Eve se aleja del valor esperado por
azar. Los propios autores reportan que la forma de cifrado que emerge de ese objetivo, sin
ninguna presión adicional, generalmente \emph{no} es un XOR bit a bit sino una función continua
que difunde cada bit de entrada sobre varios elementos del texto cifrado; también reportan que el
entrenamiento adversarial no siempre converge (hasta un tercio de las corridas fallan según sus
propios experimentos), sin diagnosticar la causa a nivel de pesos individuales. Nuestro trabajo
usa una variante de su esquema con un término auxiliar explícito que sí empuja el cifrado hacia
la forma XOR (Sección~\ref{sec:metodo}) —precisamente para tener un blanco con semántica
conocida contra el cual comparar— y retoma el problema abierto de diagnosticar por qué la
convergencia a ese blanco no siempre es completa.

\textbf{RASP y Tracr.} \citet{weiss2021thinking} proponen RASP como modelo computacional para
razonar sobre lo que un transformer puede o no expresar en función de su número de capas y
cabezas de atención. \citet{lindner2023tracr} construyen sobre RASP un compilador
(\texttt{tracr}) que traduce un programa RASP a los pesos de un transformer real, produciendo
un modelo con estructura interna completamente conocida. Los autores proponen explícitamente
usar estos modelos compilados como banco de pruebas (``laboratorio'') para técnicas de
interpretabilidad, en lugar de como objeto de estudio en sí mismos —es exactamente ese marco el
que extendemos aquí, usando el vocabulario del modelo compilado como fuente de hipótesis para
sondear una red externa.

\textbf{Sondeo lineal (probing).} La idea de entrenar clasificadores lineales sobre
representaciones internas de una red para evaluar si una propiedad está linealmente codificada
se remonta a \citet{alain2016understanding}. Usamos esa misma herramienta, pero con una
diferencia relevante: en lugar de sondear por una propiedad genérica del dominio, sondeamos
específicamente por las variables que un modelo con semántica conocida (Tracr) ya demostró que
son necesarias y suficientes para resolver la tarea.

\textbf{Interpretabilidad mecanística.} \citet{olah2020zoom} argumentan que los mecanismos
internos de una red pueden entenderse estudiando las conexiones entre neuronas como circuitos.
\citet{olah2022mechanistic} discuten la importancia de que la base de variables que usa una red
para representar un problema coincida con una base interpretable por humanos —la
interpretabilidad, argumentan, no es solo una propiedad del método usado sino también del
modelo y de si su representación natural es o no legible. Nuestro resultado puede leerse como
una instancia concreta de ese argumento: la red Alice original solo tiene una base
interpretable (una codificación lineal limpia por bit) en una fracción de sus posiciones de
salida, y el resto queda opaco hasta que se corrige un cuello de botella de capacidad.

\section{Método}
\label{sec:metodo}

\subsection{Esquema ANC}

Usamos una versión simplificada del esquema de tres redes de Abadi y Andersen. La arquitectura
original combina una capa totalmente conectada con varias capas convolucionales 1D
(``mix \& transform''); en cambio, tanto Alice como Bob aquí son un MLP de una sola capa
oculta, elegido deliberadamente así porque el objeto de estudio de este trabajo es esa capa
oculta. Con longitud de mensaje $N$, Alice recibe la concatenación de mensaje y clave ($2N$
bits), la pasa por una capa oculta totalmente conectada de tamaño \texttt{hidden\_size} con
activación ReLU, y produce un texto cifrado de $N$ bits mediante una capa de salida con
activación sigmoide. Bob tiene la misma arquitectura, recibiendo cifrado y clave. Eve, sin
acceso a la clave, recibe solo el cifrado. Por defecto \texttt{hidden\_size} $= 2N$ (el mismo
tamaño que la entrada concatenada); en la Sección~\ref{sec:resultados} tratamos
\texttt{hidden\_size} como un hiperparámetro independiente para aislar su efecto.

El objetivo de entrenamiento también es una variante del de Abadi y Andersen, que usan una
penalización cuadrática sobre qué tan lejos está el error de Eve de lo esperado por azar y no
incluyen ningún término de XOR explícito. Nuestra variante combina tres términos:
\begin{equation}
  \mathcal{L}_{\text{Alice}} = \mathcal{L}_{\text{Bob}} - \alpha\,\mathcal{L}_{\text{Eve}} + \beta\,\mathcal{L}_{\text{XOR}}
\end{equation}
donde $\mathcal{L}_{\text{Bob}}$ y $\mathcal{L}_{\text{Eve}}$ son la entropía cruzada binaria
entre la predicción de cada decodificador y el mensaje original, y $\mathcal{L}_{\text{XOR}}$
es la entropía cruzada entre el texto cifrado producido y la referencia XOR explícita
bit a bit —el término que garantiza que exista un blanco conocido contra el cual comparar,
ausente en la formulación original—. $\alpha$ pondera cuánto se premia confundir a Eve; $\beta$
pondera cuánto se empuja el cifrado hacia la forma XOR específicamente. Cada paso de
entrenamiento actualiza primero Bob, luego Alice, y finalmente Eve durante varias
sub-iteraciones (6 por defecto) para mantener
a Eve cerca de su óptimo dado el Alice actual, siguiendo el esquema de entrenamiento adversarial
estándar. Usamos Adam \citep{kingma2015adam} con tasa de aprendizaje $10^{-3}$ para las tres
redes.

\subsection{Modelo de referencia compilado con Tracr}

Implementamos el cifrado XOR como un programa RASP: la entrada es una única secuencia de largo
$2N$, con los bits del mensaje en las posiciones $[0, N)$ y los de la clave en $[N, 2N)$. Para
alinear cada bit de mensaje con su bit de clave correspondiente se usa un mecanismo de atención
expresado como \texttt{Select}/\texttt{Aggregate} sobre los índices de posición
(\texttt{key\_pos == query\_pos + N}), que ``trae'' categóricamente el valor del bit de clave a
la posición del bit de mensaje; un \texttt{SequenceMap} calcula entonces el XOR bit a bit entre
ambos valores. Compilado con \texttt{tracr.compiler.compiling.compile\_rasp\_to\_model}, el
resultado es un transformer mínimo: una capa, una cabeza de atención, \texttt{key\_size=10},
\texttt{mlp\_hidden\_size=4}, con un residual stream de dimensiones categóricas etiquetadas
(\texttt{model.residual\_labels}). Dos de esas etiquetas son el vocabulario que usamos como
hipótesis en el resto del trabajo: \texttt{aggregate\_2:\{0,1\}} (el bit de clave traído a la
posición del mensaje) y \texttt{sequence\_map\_1:\{0,1\}} (el bit de salida XOR en esa
posición). Verificamos programáticamente, leyendo los logits de atención del modelo compilado
(no el código RASP), que cada una de las $N$ posiciones de mensaje efectivamente atiende a la
posición \texttt{pos+N} ($N=4$ en nuestra verificación). Como el mecanismo de \texttt{Select} de
este programa RASP depende únicamente de índices de posición y no de los valores de los bits,
el patrón de atención es idéntico para cualquier mensaje y clave —una sola verificación alcanza
para confirmar el mecanismo en general, y este es legible directamente en los pesos, sin
necesidad de inspeccionar el código RASP fuente.

\subsection{Sondeo lineal guiado por el vocabulario de Tracr}
\label{sec:probing}

Dado que Alice es un MLP de una sola capa oculta y Tracr compila un transformer, la comparación
no puede hacerse circuito a circuito. En cambio, tomamos la variable \texttt{sequence\_map\_1}
(el bit XOR de salida en la posición $i$) como hipótesis y preguntamos: ¿está esa misma
información codificada de forma linealmente legible en la capa oculta de Alice? Para cada
posición de salida $i \in [0, N)$ entrenamos una sonda lineal (regresión logística, implementada
como una capa \texttt{nn.Linear} entrenada con \texttt{BCEWithLogitsLoss} y Adam, sin
dependencias externas de \emph{probing}) que predice el bit XOR de esa posición a partir de dos
posibles fuentes de features:

\begin{itemize}
\item \textbf{\texttt{input}} (control): la concatenación cruda de mensaje y clave. Como el
  valor de XOR en una posición no es linealmente separable en función de únicamente esos dos
  bits relevantes —el problema clásico de separabilidad de XOR
  \citep{minsky1969perceptrons}—, este control acota por debajo lo que puede lograr una sonda
  lineal si Alice no realizó ningún cómputo no lineal útil.
\item \textbf{\texttt{hidden}}: las activaciones post-ReLU de la única capa oculta de Alice
  (\texttt{fc1}), el análogo funcional del residual stream de Tracr.
\end{itemize}

Cada sonda se entrena con 4096 ejemplos y se evalúa sobre 1024 ejemplos held-out generados de
manera independiente.

\subsection{Diagnóstico de colapso y ablaciones}

Al observar que la accuracy de la sonda \texttt{hidden} es marcadamente bimodal entre
posiciones (cercana a 1.0 o cercana al azar, rara vez intermedia), inspeccionamos la norma de
cada fila de la matriz de pesos \texttt{fc2.weight} de Alice (la conexión capa oculta
$\rightarrow$ salida), una fila por posición de salida. Definimos una posición como
\emph{colapsada} si esa norma cae por debajo de un umbral de 4.0. El umbral se eligió post-hoc,
pero está bien justificado empíricamente: agrupando las normas de fila observadas en los 280
pares (posición, configuración) de las ablaciones de la Sección~\ref{sec:resultados}, existe un
hueco limpio entre $3.68$ (el valor más alto por debajo del umbral) y $4.25$ (el valor más bajo
por encima), sin ninguna observación en ese intervalo.

Para investigar la causa del colapso realizamos dos barridos, cada uno con 5 semillas
aleatorias fijas (0--4) y 10\,000 pasos de entrenamiento, $N=8$, tamaño de minibatch 128:

\begin{enumerate}
\item \textbf{Barrido de pesos de pérdida}: 7 combinaciones de $(\alpha, \beta)$ ---
  $(1,0.5)$, $(1,0.1)$, $(2,0.5)$, $(2,0.1)$, $(3,0.5)$, $(2,1.0)$ y $(3,1.0)$ ---
  con \texttt{hidden\_size} fijo en $2N=16$ (arquitectura original).
\item \textbf{Barrido de capacidad}: \texttt{hidden\_size} $\in \{16, 32, 64, 128\}$, con
  $(\alpha,\beta)$ fijo en la mejor combinación encontrada en el barrido anterior
  ($3.0, 1.0$), más una verificación de aislamiento cruzando \texttt{hidden\_size} $\in
  \{16,32\}$ con $(\alpha,\beta) \in \{(1.0,0.5), (3.0,1.0)\}$ para separar el efecto de la
  capacidad del efecto del balance de pérdidas.
\end{enumerate}

En cada configuración medimos el \texttt{xor\_match} promedio (fracción de bits de salida que
coinciden con la referencia XOR explícita, promediada sobre posiciones y sobre un batch de
evaluación de 4096 ejemplos) y el número de posiciones colapsadas según el criterio de norma
anterior.

\section{Resultados}
\label{sec:resultados}

\subsection{El sondeo revela heterogeneidad invisible en la métrica agregada}

Con $N=8$, $\alpha=1.0$, $\beta=0.5$ (valores por defecto) y 22\,000 pasos de entrenamiento
(semilla 0), el \texttt{xor\_match} agregado se estabiliza en $\approx 0.83$, sin llegar nunca
a 1.0. La Tabla~\ref{tab:probe} muestra la accuracy de la sonda lineal por posición para las
dos fuentes de features.

\begin{table}[h]
  \caption{Accuracy de la sonda lineal (test set) por posición de salida, $N=8$, semilla 0,
    22\,000 pasos de entrenamiento.}
  \label{tab:probe}
  \begin{tabular}{lcccccccc}
    \toprule
    Posición & 0 & 1 & 2 & 3 & 4 & 5 & 6 & 7 \\
    \midrule
    \texttt{input} (control) & 0.501 & 0.485 & 0.467 & 0.529 & 0.479 & 0.521 & 0.503 & 0.515 \\
    \texttt{hidden} (Alice)  & 1.000 & 0.500 & 0.521 & 1.000 & 0.634 & 1.000 & 1.000 & 1.000 \\
    \bottomrule
  \end{tabular}
\end{table}

El control \texttt{input} se mantiene, como se predijo, cercano al azar en las 8 posiciones —el
límite teórico de separabilidad lineal para XOR. La sonda \texttt{hidden}, en cambio, es
marcadamente bimodal: 5 de 8 posiciones (0, 3, 5, 6, 7) son legibles con accuracy perfecta,
tan limpiamente como la variable categórica \texttt{sequence\_map\_1} en el modelo compilado con
Tracr, mientras que las posiciones 1 y 2 quedan al nivel de azar y la posición 4 es intermedia.
Esta heterogeneidad es consistente aritméticamente con el \texttt{xor\_match} agregado
observado: si 5 posiciones están resueltas perfectamente y las 3 restantes al azar, el promedio
esperado es $(5 \times 1.0 + 3 \times 0.5)/8 \approx 0.8125$, casi exactamente el valor medido.
Es decir, las posiciones donde la sonda lineal lee el bit XOR con perfección son las mismas
donde Alice efectivamente convergió al cifrado correcto.

\subsection{El colapso es consistente entre semillas y se explica por la norma de los pesos de salida}

Repetimos el entrenamiento con 5 semillas distintas (10\,000 pasos, $\alpha=1.0$,
$\beta=0.5$) y medimos, para cada una, el \texttt{xor\_match} de cada posición individual y la
norma de la fila correspondiente de \texttt{fc2.weight} —40 pares (norma, \texttt{xor\_match})
en total. La correlación de Pearson entre ambas cantidades es $r=0.97$. La relación no es un
escalón perfecto entre azar y perfección: hay casos intermedios genuinos en ambos extremos del
umbral (p.~ej.\ norma $=3.68 \Rightarrow \texttt{xor\_match}=0.872$; norma $=6.05 \Rightarrow
\texttt{xor\_match}=0.875$), pero el patrón es inequívoco —dentro del grupo con norma por debajo
del umbral de $4.0$ definido en la sección anterior, el \texttt{xor\_match} por posición varía entre
$0.481$ y $0.872$ (media $0.559$); dentro del grupo con norma por encima del umbral, varía entre
$0.875$ y $1.000$ (media $0.991$), sin superposición entre ambos rangos. Ninguna semilla
presentó unidades ReLU muertas en la capa oculta (0/16 en las 5 semillas), descartando esa
hipótesis alternativa: el colapso está específicamente en la conexión capa oculta
$\rightarrow$ salida de la posición afectada, no en la representación oculta compartida. Qué
posición específica colapsa varía con la semilla (no hay una posición estructuralmente
privilegiada), pero el mecanismo es siempre el mismo.

Interpretamos este colapso como un óptimo local estable del objetivo minimax de tres redes:
cuando el peso de salida de un bit colapsa hacia cero, el sigmoide converge a una salida
aproximadamente constante $\approx 0.5$ para ese bit, independiente del input. En ese régimen,
Bob acierta ese bit solo $\approx 50\%$ del tiempo (penalizando $\mathcal{L}_{\text{Bob}}$),
pero Eve también queda en $\approx 50\%$ (beneficiando el término $-\alpha \mathcal{L}_{\text{Eve}}$)
—``rendirse'' en un bit específico intercambia precisión de Bob por confusión de Eve, y como el
target de ese bit es aleatorio 50/50, el gradiente que debería revertir el colapso se cancela en
promedio, dejando el punto como un atractor estable.

\subsection{El balance de la función de pérdida modula el colapso pero no lo elimina}

\begin{table}[h]
  \caption{Barrido de pesos de pérdida $(\alpha,\beta)$, $N=8$, \texttt{hidden\_size}$=16$,
    10\,000 pasos, promedio $\pm$ desviación estándar sobre 5 semillas.}
  \label{tab:alphabeta}
  \begin{tabular}{lcc}
    \toprule
    Configuración & \texttt{xor\_match} & Posiciones colapsadas (/8) \\
    \midrule
    $\alpha{=}1.0,\ \beta{=}0.5$ (original) & $0.850 \pm 0.059$ & $2.60 \pm 1.02$ \\
    $\alpha{=}1.0,\ \beta{=}0.1$            & $0.500 \pm 0.002$ & $8.00 \pm 0.00$ \\
    $\alpha{=}2.0,\ \beta{=}0.5$            & $0.915 \pm 0.032$ & $1.60 \pm 0.49$ \\
    $\alpha{=}2.0,\ \beta{=}0.1$            & $0.504 \pm 0.007$ & $8.00 \pm 0.00$ \\
    $\alpha{=}3.0,\ \beta{=}0.5$            & $0.871 \pm 0.045$ & $2.40 \pm 0.80$ \\
    $\alpha{=}2.0,\ \beta{=}1.0$            & $0.916 \pm 0.022$ & $1.60 \pm 0.49$ \\
    $\alpha{=}3.0,\ \beta{=}1.0$            & $0.921 \pm 0.014$ & $1.60 \pm 0.49$ \\
    \bottomrule
  \end{tabular}
\end{table}

La Tabla~\ref{tab:alphabeta} resume el barrido de pesos de pérdida. Bajar $\beta$ no reduce el
colapso: lo generaliza a las 8 posiciones, pero por un mecanismo distinto. Con
$\beta=0.1$, \texttt{bob\_acc} sube a $\approx 0.97$--$1.0$ mientras \texttt{xor\_match} cae al
azar en las 8 posiciones —Alice deja de intentar converger a la forma XOR específicamente, y en
cambio encuentra otra función (no-XOR) para comunicarse con Bob. $\beta$ es entonces lo que
fuerza la convergencia a la forma XOR, no un término prescindible.

Subir $\alpha$ manteniendo $\beta=0.5$ sí ayuda: con $\alpha=2.0$ el promedio de posiciones
colapsadas baja de $2.60$ a $1.60$ y \texttt{xor\_match} sube de $0.850$ a $0.915$. Sin embargo,
subir $\alpha$ aún más ($3.0$) sin acompañarlo de un aumento de $\beta$ es \emph{peor} que
$\alpha=2.0$ ($2.40$ vs. $1.60$ posiciones colapsadas) —el efecto no es monótono. Las tres
configuraciones con mejor desempeño ($\alpha{=}2/\beta{=}0.5$, $\alpha{=}2/\beta{=}1.0$,
$\alpha{=}3/\beta{=}1.0$) convergen al mismo piso de $1.60 \pm 0.49$ posiciones colapsadas en
promedio, sugiriendo que el ajuste de estos dos hiperparámetros por sí solo, dentro del rango
explorado, no elimina el fenómeno.

\subsection{El colapso es principalmente un cuello de botella de capacidad}

\begin{table}[h]
  \caption{Barrido de \texttt{hidden\_size}, $N=8$, $\alpha=3.0$, $\beta=1.0$, 10\,000 pasos,
    promedio $\pm$ desviación estándar sobre 5 semillas.}
  \label{tab:capacity}
  \begin{tabular}{lcc}
    \toprule
    \texttt{hidden\_size} & \texttt{xor\_match} & Posiciones colapsadas (/8) \\
    \midrule
    16 ($=2N$, original) & $0.921 \pm 0.014$ & $1.60 \pm 0.49$ \\
    32 ($=4N$)            & $1.000 \pm 0.000$ & $0.00 \pm 0.00$ \\
    64 ($=8N$)            & $1.000 \pm 0.000$ & $0.00 \pm 0.00$ \\
    128 ($=16N$)          & $1.000 \pm 0.000$ & $0.00 \pm 0.00$ \\
    \bottomrule
  \end{tabular}
\end{table}

La Tabla~\ref{tab:capacity} muestra el resultado más contundente del trabajo: duplicar el ancho
de la capa oculta ($16 \to 32$) elimina el colapso por completo, en las 5 semillas, sin
excepción. Ampliar aún más (64, 128) no aporta ninguna mejora adicional —el techo ya se alcanza
en $4N$.

Para aislar si el efecto es puramente de capacidad o requiere también el ajuste de pérdidas de
la sección anterior, cruzamos \texttt{hidden\_size} $\in \{16, 32\}$ con
$(\alpha,\beta) \in \{(1.0,0.5)_{\text{original}}, (3.0,1.0)_{\text{ajustado}}\}$
(Tabla~\ref{tab:isolation}).

\begin{table}[h]
  \caption{Aislamiento del efecto de capacidad vs.\ balance de pérdidas, $N=8$, 10\,000
    pasos, promedio $\pm$ desviación estándar sobre 5 semillas.}
  \label{tab:isolation}
  \begin{tabular}{llcc}
    \toprule
    \texttt{hidden\_size} & $(\alpha,\beta)$ & \texttt{xor\_match} & Colapsadas (/8) \\
    \midrule
    16 & original ($1.0,0.5$) & $0.850 \pm 0.059$ & $2.60 \pm 1.02$ \\
    32 & original ($1.0,0.5$) & $0.958 \pm 0.035$ & $0.40 \pm 0.49$ \\
    16 & ajustado ($3.0,1.0$) & $0.921 \pm 0.014$ & $1.60 \pm 0.49$ \\
    32 & ajustado ($3.0,1.0$) & $1.000 \pm 0.000$ & $0.00 \pm 0.00$ \\
    \bottomrule
  \end{tabular}
\end{table}

Los dos efectos son reales, independientes, y se combinan. Más capacidad por sí sola
(\texttt{hidden\_size}$=32$ con los pesos de pérdida \emph{originales}) ya reduce el colapso de
$2.60$ a $0.40$ posiciones y sube \texttt{xor\_match} de $0.850$ a $0.958$ —la capacidad explica
la mayor parte del fenómeno. Pero ninguno de los dos factores por sí solo alcanza el resultado
perfecto: se necesita capacidad suficiente ($\texttt{hidden\_size} \geq 32$, el doble de la
arquitectura original) \emph{y} el reajuste de $(\alpha,\beta)$ juntos para llegar a
$\texttt{xor\_match}=1.000$ con cero posiciones colapsadas en las 5 semillas sin excepción.

\section{Discusión y conclusiones}

\subsection{El vocabulario de Tracr como generador de hipótesis, no como comparación circuito a circuito}

El resultado central de este trabajo no depende de que Alice (un MLP) y el modelo compilado
(un transformer) compartan estructura de circuito —no la comparten, y no podrían compararse
directamente peso a peso. Lo que transferimos de Tracr a AliceNet no es una arquitectura sino
un \emph{vocabulario de variables}: la noción de que existe, para cada posición de salida, un
bit categórico (el XOR de esa posición) que en el modelo con semántica conocida es legible al
100\% de forma lineal. Ese vocabulario se convierte en una hipótesis falsable sobre la red
opaca (\S\ref{sec:probing}), y resulta que la hipótesis es cierta solo parcialmente —lo cual es
en sí mismo el hallazgo. Esto es consistente con el argumento de \citet{olah2022mechanistic}
de que la interpretabilidad de un modelo depende de si su base de representación natural
coincide con una base interpretable por humanos: aquí mostramos una forma operacional y
falsable de poner a prueba esa coincidencia, usando un modelo con base conocida por
construcción como fuente de la hipótesis.

\subsection{De síntoma agregado a diagnóstico mecánico}

Sin el sondeo por posición, el único síntoma visible habría sido que \texttt{xor\_match} se
estanca en $\approx 0.83$--$0.92$ sin llegar a 1.0, algo que admite muchas explicaciones
posibles (ruido de entrenamiento, número insuficiente de pasos, mala elección de
hiperparámetros del optimizador). El sondeo lineal por posición, guiado por el vocabulario de
Tracr, redujo el espacio de hipótesis a una específica y verificable: hay una subred de
posiciones que efectivamente colapsó. Inspeccionar la norma de \texttt{fc2.weight} confirmó el
mecanismo exacto (colapso de la conexión capa oculta $\to$ salida, no de la capa oculta en sí),
y la ablación de capacidad confirmó la causa raíz (capacidad insuficiente de la capa oculta
compartida). Este encadenamiento —síntoma agregado, hipótesis guiada por vocabulario
interpretable, sondeo, inspección de pesos, ablación causal— es reproducible como metodología
general para diagnosticar por qué una red adversarial no converge al comportamiento esperado,
más allá del caso particular de XOR.

\subsection{Limitaciones}

Este trabajo tiene limitaciones que acotan el alcance de sus conclusiones. Primero, la tarea
(XOR bit a bit, $N=8$) es deliberadamente simple; no está claro que el mismo patrón de colapso
por posición aparezca, o tenga la misma causa, en cifrados más complejos o en tareas fuera del
dominio criptográfico. Segundo, el umbral de colapso (norma $<4.0$) se eligió post-hoc a partir
del hueco observado empíricamente en la distribución agregada de normas
(\S\ref{sec:metodo}), y no de un criterio independiente del dato; aunque la separación de norma
es limpia, el \texttt{xor\_match} asociado a cada lado del umbral no lo es igual de tajantemente
—existen posiciones con norma justo por debajo del umbral cuyo \texttt{xor\_match} ya es
sustancialmente mayor que el azar (\S\ref{sec:resultados})—, y un umbral fijo podría no
generalizar a otras escalas de $N$ o de \texttt{hidden\_size}. Tercero,
el sondeo lineal establece legibilidad, no necesidad causal: no realizamos intervenciones
(activation patching) dentro de la propia red Alice para confirmar que la unidad de información
identificada es efectivamente usada causalmente por el resto de la red, solo inferimos
causalidad indirectamente a través de la ablación de arquitectura (reentrenamientos completos
con distinta capacidad). Cuarto, el número de semillas (5) por configuración, aunque suficiente
para mostrar el patrón cualitativo de forma consistente, es modesto para una caracterización
estadística fina de los promedios reportados.

\subsection{Conclusiones}

Mostramos un uso concreto del vocabulario de variables expuesto por un modelo compilado con
Tracr como generador de hipótesis para un sondeo lineal aplicado a una red de criptografía
neuronal adversarial entrenada de forma opaca. Ese sondeo reveló una heterogeneidad de
convergencia por posición de salida invisible en la métrica agregada habitual, y permitió
diagnosticar mecánicamente su causa: un colapso de la norma de los pesos de salida específico
por posición, consistente con un óptimo local del objetivo minimax de tres redes. Mediante una
ablación controlada mostramos que este colapso es principalmente un cuello de botella de
capacidad de la arquitectura original de Abadi y Andersen —duplicar el ancho de la capa oculta
lo reduce drásticamente incluso sin tocar los pesos de la función de pérdida, y combinado con un
reajuste de esos pesos lo elimina por completo— y no una limitación fundamental del esquema de
entrenamiento adversarial en sí. El resultado ilustra que el valor de un modelo con ground-truth
compilado no se limita a validar técnicas de interpretabilidad en abstracto: su vocabulario de
variables puede transferirse directamente como hipótesis de diagnóstico sobre una red opaca no
relacionada arquitectónicamente, y ese diagnóstico puede guiar una corrección de diseño
concreta y verificable.

\section*{Declaración de ética}

Este trabajo no involucra datos personales, sujetos humanos, ni aplicaciones de doble uso más
allá de las ya inherentes a la investigación en criptografía neuronal, un dominio establecido en
la literatura de aprendizaje automático desde \citet{abadi2016learning}. Todos los experimentos
se realizaron sobre datos sintéticos (mensajes y claves generados aleatoriamente) y modelos de
juguete a pequeña escala, sin riesgos de seguridad prácticos.

\begin{acks}
Este proyecto está financiado por InsaneMonkey SpA --- MonkeysLab Unhinged Research
(\url{https://www.monkeyslab.cl}).
\end{acks}

\bibliographystyle{ACM-Reference-Format}
\bibliography{referencias}

\end{document}
